% !TEX root = ../matan.tex

\section{Дифференцируемость функций нескольких переменных. Непрерывность дифференцируемой функции.}

Пусть $\Omega\subset\R^d$ открыто, $x\in\Omega$, $f\colon\Omega\to\R^n$.

\begin{Definition}[required]{Дифференцируемость}{}
	Функция $f$ называется \emph{дифференцируемой} в точке $x$, если существует линейный оператор $A=d_xf\colon\R^d\to\R^n$ такой, что при $\xi\to x$, $\xi\in\Omega$,
	\[
	f(\xi)=f(x)+A(\xi-x)+o(|\xi-x|),
	\]
	то есть
	\[
	\frac{|f(\xi)-f(x)-A(\xi-x)|}{|\xi-x|}\xrightarrow[\xi\to x]{}0.
	\]
	Оператор $A$ называется \emph{дифференциалом (производной)} функции $f$ в точке $x$ и обозначается $d_xf$ (также $f'(x)$, $Df(x)$). Геометрически $\xi\mapsto f(x)+A(\xi-x)$ --- наилучшее аффинное приближение $f$ вблизи $x$.
\end{Definition}

\begin{Remark}{Связь с одномерным случаем}{}
	При $d=n=1$ линейный оператор $A\colon\R\to\R$ есть умножение на число, $A(h)=f'(x)\,h$, и определение превращается в обычное:
	\[
	f(\xi)=f(x)+f'(x)(\xi-x)+o(|\xi-x|).
	\]
\end{Remark}

\begin{Statement}{Единственность дифференциала}{}
	Если линейный оператор $A$ из определения существует, то он определён однозначно.
\end{Statement}

\begin{proof}
	% [восстановлено]
	Пусть операторы $A$ и $B$ оба удовлетворяют определению. Вычитая два разложения для $f(\xi)$, получаем
	\[
	(A-B)(\xi-x)=o(|\xi-x|)\qquad(\xi\to x).
	\]
	Зафиксируем произвольный вектор $v\ne0$ и подставим $\xi=x+tv$ с $t\to0$, $t\ne0$ (при малых $t$ точка $\xi\in\Omega$, так как $\Omega$ открыто). Тогда $\xi-x=tv$, $|\xi-x|=|t|\,|v|$, и по линейности
	\[
	t\,(A-B)v=(A-B)(tv)=o(|t|),
	\qquad\text{откуда}\qquad
	(A-B)v=\frac{o(|t|)}{t}\xrightarrow[t\to0]{}0.
	\]
	Левая часть от $t$ не зависит, значит $(A-B)v=0$. Так как вектор $v$ произволен, $A-B=0$, то есть $A=B$.
\end{proof}

\begin{Theorem}[required]{Дифференцируемость влечёт непрерывность}{}
	Если $f\colon\Omega\to\R^n$ дифференцируема в точке $x$, то $f$ непрерывна в точке $x$.
\end{Theorem}

\begin{proof}
	Любой линейный оператор $A\colon\R^d\to\R^n$ ограничен: существует $C\ge0$ такое, что $|Ah|\le C|h|$ для всех $h\in\R^d$. Действительно, если $A=(A_{ij})$ --- его матрица, то по неравенству Коши--Буняковского $|Ah|\le C|h|$ с $C=\sqrt{\sum_{i,j}A_{ij}^2}$.

	Из определения дифференцируемости $f(\xi)-f(x)=d_xf(\xi-x)+o(|\xi-x|)$, поэтому
	\[
	|f(\xi)-f(x)|\le |d_xf(\xi-x)|+o(|\xi-x|)\le C\,|\xi-x|+o(|\xi-x|)\xrightarrow[\xi\to x]{}0.
	\]
	Значит, $f(\xi)\to f(x)$ при $\xi\to x$, то есть $f$ непрерывна в точке $x$.
\end{proof}

\begin{Theorem}{Покомпонентная дифференцируемость}{}
	Пусть $f=(f_1,\dots,f_n)\colon\Omega\subset\R^d\to\R^n$. Тогда $f$ дифференцируема в точке $x$ в том и только том случае, когда каждая координатная функция $f_k\colon\Omega\to\R$ дифференцируема в точке $x$; при этом $k$-й компонентой дифференциала служит дифференциал $f_k$:
	\[
	d_xf=\begin{pmatrix} d_xf_1\\ \vdots\\ d_xf_n\end{pmatrix},
	\qquad\text{то есть}\quad \bigl(d_xf(h)\bigr)_k=d_xf_k(h).
	\]
\end{Theorem}

\begin{proof}
	Пусть $A\colon\R^d\to\R^n$ --- линейный оператор, $A_k$ --- его $k$-я компонента (линейный функционал $h\mapsto(Ah)_k$). Рассмотрим вектор-остаток
	\[
	\rho(\xi)=\frac{f(\xi)-f(x)-A(\xi-x)}{|\xi-x|}\in\R^n,\qquad
	\rho_k(\xi)=\frac{f_k(\xi)-f_k(x)-A_k(\xi-x)}{|\xi-x|}.
	\]
	По доказанному в вопросе~37 двустороннему неравенству $|\rho_k(\xi)|\le|\rho(\xi)|\le\sqrt n\,\max_k|\rho_k(\xi)|$ покоординатная сходимость $\rho(\xi)\to0$ равносильна сходимости $\rho_k(\xi)\to0$ для всех $k$. Условие $\rho(\xi)\to0$ означает дифференцируемость $f$ с дифференциалом $A$, а условие $\rho_k(\xi)\to0$ при всех $k$ --- дифференцируемость каждой $f_k$ с дифференциалом $A_k$. Отсюда эквивалентность и указанная формула для $d_xf$.
\end{proof}
